
% unit: noether.p17.s11.title
\subsection*{\S~11. Relations avec les systèmes d'équations différentielles et leurs intégrales.}

% unit: noether.p17.s11.par.001
En vertu du théorème de finitude, on peut associer à tout module d'expressions différentielles un système composé d'un nombre fini d'expressions différentielles comme base, de telle sorte que l'ensemble des intégrales simultanées de toutes les équations différentielles obtenues en annulant une expression différentielle du module coïncide avec l'ensemble des intégrales du système fini d'équations différentielles que l'on obtient en annulant les expressions différentielles de base. C'est ainsi que se donne le rattachement à la théorie des systèmes d'équations différentielles partielles linéaires; nous poursuivons encore un peu ce rattachement dans le cas où le groupe quotient du module est \og{} fini \fg{}.

% unit: noether.p17.s11.def.001
\textbf{Définition.} \emph{Un groupe quotient s'appelle fini s'il existe un nombre $\mu$, l'ordre du groupe, tel qu'entre $\mu+1$ classes résiduelles quelconques il existe une relation linéaire à coefficients dans $P$, tandis qu'il existe au moins un système de $\mu$ classes résiduelles entre lesquelles aucune relation n'existe, système que nous appellerons une base linéaire du groupe quotient. Dans l'autre cas, le groupe quotient s'appelle infini.}

% unit: noether.p17.s11.par.002
La somme de deux groupes quotients finis a un ordre fini, à savoir la somme de leurs ordres. La somme de deux groupes infinis, ou d'un groupe fini et d'un groupe infini, est un groupe infini. Des exemples de groupes finis sont fournis par tous les groupes quotients de modules à une variable; mais, même pour plusieurs variables, l'ordre peut devenir fini, par exemple pour le module de base $\xi_1^2$, $\xi_2^2$, où toutes les classes résiduelles peuvent être composées linéairement à partir de celles de $1,\xi_1,\xi_2,\xi_1\xi_2$. Pour des exemples de groupes infinis, voir \S~12.

% unit: noether.p17.s11.par.003
On a maintenant le théorème suivant.

% unit: noether.p17.s11.thm.011
\textbf{Théorème XI.} \emph{Si un module possède un groupe quotient fini, alors tout système d'équations différentielles qui naît de l'annulation d'une base du module admet exactement autant d'intégrales linéairement indépendantes que l'indique l'ordre de son groupe quotient. Réciproquement, pour un système donné d'un nombre fini d'équations différentielles partielles linéaires qui ne possèdent qu'un nombre fini d'intégrales linéairement indépendantes, ce nombre est égal à l'ordre du groupe quotient du module engendré par les expressions différentielles.}

% unit: noether.p17.s11.par.004
Le domaine de rationalité $P$ est ici constitué de fonctions analytiques de $n$ variables $x_1,\ldots,x_n$, dont les singularités éventuelles, à l'intérieur d'un certain domaine $n$-dimensionnel, forment tout au plus des variétés de dimensions inférieures, de sorte qu'il existe un nombre arbitraire de points ayant des voisinages réguliers $n$-dimensionnels.

% unit: noether.p17.s11.par.005
Soit maintenant $\mM$ un module à groupe quotient fini; nous montrons alors qu'il existe une base linéaire formée de produits de puissances qui possède la propriété suivante: \emph{dans l'expression d'une classe résiduelle quelconque au moyen de cette base, seules des puissances d'un nombre fini de grandeurs différentes de $P$ peuvent apparaître au dénominateur}.

% unit: noether.p17.s11.par.006
À cette fin, nous ordonnons lexicographiquement tous les produits de puissances $\xi_1^{a_1}\cdots\xi_n^{a_n}$ et prenons comme représentants $Q_1,\ldots,Q_m$ des classes résiduelles tous ceux qui sont linéairement indépendants modulo $\mM$ de ceux qui les précèdent. Par hypothèse, il n'y en a qu'un nombre fini. Tous les autres peuvent s'exprimer linéairement par ceux-ci.

% unit: noether.p17.s11.par.007
Soit $\mathfrak S$ le système de tous les produits de puissances non admis dans la base linéaire. Alors, d'après la première partie de la démonstration du théorème III, $\mathfrak S$ possède un sous-système fini $\mathfrak T$ de tels produits de puissances, tel que tout produit de puissances de $\mathfrak S$ soit divisible par au moins un produit de $\mathfrak T$. Soient $P_1,\ldots,P_\rho$ les produits de puissances de $\mathfrak T$, et soient $M_1\equiv0,\ldots,M_\rho\equiv0\;(\mM)$ les relations au moyen desquelles $P_1,\ldots,P_\rho$ sont chaque fois exprimés, modulo $\mM$, par une combinaison linéaire de produits de puissances inférieurs; par exemple
% unit: noether.p17.s11.eq.001
\[
 M_i=b_iP_i+\text{termes inférieurs}.
\]
Si maintenant $P$ est un produit de puissances quelconque de $\mathfrak S$, donc de la forme
% unit: noether.p17.s11.eq.002
\[
 \xi_1^{h_1}\cdots\xi_n^{h_n}P_i,
\]
il vient
% unit: noether.p17.s11.eq.003
\[
\begin{aligned}
 \xi_1^{h_1}\cdots\xi_n^{h_n}M_i
 &=b_i\xi_1^{h_1}\cdots\xi_n^{h_n}P_i+\text{termes inférieurs}\\
 &=b_iP+\text{termes inférieurs}.
\end{aligned}
\]
Admettons maintenant comme démontré que, pour tous les produits de puissances inférieurs à $P$, seules des puissances de $b_1,\ldots,b_\rho$ apparaissent au dénominateur, tandis que les numérateurs sont formés à partir des coefficients des $M_i$ par addition, multiplication et différentiation; alors cela vaut aussi pour $P$ lui-même. Or cette hypothèse est permise, puisqu'elle est satisfaite pour $P_1$, le premier élément de $\mathfrak S$ et de $\mathfrak T$.

% unit: noether.p17.s11.par.008
D'après le théorème III, $M_1,\ldots,M_\rho$ est en outre une base du module $\mM$. Considérons maintenant un système de valeurs $\beta_1,\ldots,\beta_n$ de $x_1,\ldots,x_n$ pour lequel $b_1\ne0,\ldots,b_\rho\ne0$ et les autres coefficients des $M_i$ sont réguliers, c'est-à-dire un point régulier du système d'équations différentielles $M_1=0,\ldots,M_\rho=0$. Aux produits de puissances $Q_i=\xi_1^{h_{i1}}\cdots\xi_n^{h_{in}}$ $(i=1,\ldots,m)$ correspondent alors les dérivées partielles
% unit: noether.p17.s11.eq.004
\[
 \frac{\partial^{h_{i1}+\cdots+h_{in}}}{\partial x_1^{h_{i1}}\cdots\partial x_n^{h_{in}}}y,
\]
pour lesquelles nous prescrivons des valeurs constantes arbitraires $(\gamma_1,\ldots,\gamma_m)$ au point $(\beta_1,\ldots,\beta_n)$. Alors les équations différentielles $M_i=0$ déterminent de façon unique, comme valeurs finies en ce point, les valeurs de toutes les dérivées supérieures de $y$, puisque seuls les $b_i$ apparaissent au dénominateur; et l'on sait que cela implique l'existence de $m$ intégrales analytiques linéairement indépendantes.

% unit: noether.p17.s11.par.009
D'autre part, sous nos hypothèses, il ne peut exister plus de $m$ intégrales linéairement indépendantes; sinon, en un point régulier, on pourrait prescrire arbitrairement les valeurs de $m+1$ dérivées convenablement choisies, contrairement au fait qu'entre $m+1$ produits de puissances quelconques doivent exister des dépendances linéaires modulo $\mM$.

% unit: noether.p17.s11.par.010
Réciproquement, si un module possède un groupe quotient infini, on montre par la même voie qu'il existe une infinité d'intégrales linéairement indépendantes; dans le cas où il n'y a que $m$ intégrales linéairement indépendantes, le groupe quotient est donc fini, et son ordre est, d'après ce qui précède, égal à $m$.

% unit: noether.p17.s11.par.011
Pour terminer ce paragraphe, nous montrons que, pour deux modules d'expressions différentielles, le concept d'isomorphie, ou ce qui revient au même celui de même espèce, a pour les intégrales des systèmes d'équations différentielles associés la même signification que le concept d'espèce utilisé par Poincaré pour une variable (cf. Blumberg, loc. cit., appendice I).

% unit: noether.p17.s11.thm.012
\textbf{Théorème XII.} \emph{Si deux modules $\mM$ et $\mN$ d'expressions différentielles sont de même espèce, alors il existe deux expressions différentielles $P$ et $Q$ telles que $P(y)$ et $Q(z)$ parcourent respectivement toutes les intégrales de $\mN$ et de $\mM$, lorsque $y$ parcourt toutes les intégrales de $\mM$ et que $z$ parcourt toutes les intégrales de $\mN$.}

% unit: noether.p17.s11.par.012
Par hypothèse, en effet (d'après \S~9), les relations suivantes existent pour deux polynômes $P$ et $Q$ convenablement choisis:
% unit: noether.p17.s11.eq.005
\[
 MQ\equiv0(\mN),\qquad PQ\equiv1(\mN),
\]
% unit: noether.p17.s11.eq.006
\[
 NP\equiv0(\mM),\qquad QP\equiv1(\mM).
\]
Ainsi, puisque $NP(y)=0$, la fonction $P(y)$ est une intégrale de $\mN$. De même, $Q(z)$ est une intégrale de $\mM$. Puisque $PQ(z)=z$ et $QP(y)=y$, $P(y)$ parcourt toutes les intégrales de $\mN$, et $Q(z)$ toutes les intégrales de $\mM$. De cette manière, à tout $y$ non nul correspond aussi un $z$ non nul, et réciproquement.

% unit: noether.p17.s12.title
\subsection*{\S~12. Exemples.}

% unit: noether.p17.s12.par.001
1. Comme premier exemple, considérons une équation différentielle linéaire ordinaire du second ordre dont les coefficients soient représentés comme fonctions symétriques des intégrales. Le domaine de rationalité $P$ consiste ici en toutes les fonctions rationnelles des intégrales et de leurs dérivées; nous nous restreignons à un voisinage suffisamment petit d'un point régulier de l'équation différentielle. Nous définissons le module $\mM$ comme l'ensemble de toutes les expressions différentielles qui ont $y_1$ et $y_2$ pour intégrales, et pour lesquelles nous écrivons de nouveau des polynômes en une indéterminée $\xi$. La base du module $\mM$ est alors à un seul élément, et elle est donnée par le polynôme
% unit: noether.p17.s12.eq.001
\[
 M=
 \begin{vmatrix} y_1&y_2\\ y_1'&y_2'\end{vmatrix}\xi^2
 -\begin{vmatrix} y_1&y_2\\ y_1''&y_2''\end{vmatrix}\xi
 +\begin{vmatrix} y_1'&y_2'\\ y_1''&y_2''\end{vmatrix},
\]
qui est déterminé de façon unique à un facteur de $P$ près. On a alors
% unit: noether.p17.s12.eq.041
\begin{equation}
 \mM=[\mM_1,\mM_2],
\tag{41}
\end{equation}
où $\mM_i$ est constitué de toutes les expressions différentielles possédant l'intégrale $y_i$, et dont la base est déterminée à un facteur de $P$ près par
% unit: noether.p17.s12.eq.002
\[
 M_i=y_i\xi-y_i'\qquad(i=1,2).
\]
En effet, $\mM$ est divisible par $\mM_1$ et $\mM_2$, puisqu'il s'annule pour les intégrales $y_1$ et $y_2$, donc il est divisible par $[\mM_1,\mM_2]$; et, comme l'ordre du polynôme générateur de $[\mM_1,\mM_2]$ doit être au moins $2$, on a $\mM=[\mM_1,\mM_2]$. De plus, $\mM_1$ et $\mM_2$ sont premiers entre eux, car
% unit: noether.p17.s12.eq.042
\begin{equation}
 \frac{y_2}{D}(y_1\xi-y_1')-\frac{y_1}{D}(y_2\xi-y_2')=1,
 \qquad
 D=D(y)=\begin{vmatrix}y_1&y_2\\ y_1'&y_2'\end{vmatrix}.
\tag{42}
\end{equation}
Ils sont en outre des modules premiers, puisque leur groupe quotient a l'ordre 1.

% unit: noether.p17.s12.par.002
Mais, outre $y_1$ et $y_2$, $\mM$ possède aussi toutes les intégrales
% unit: noether.p17.s12.eq.003
\[
 z_1=\alpha_{11}y_1+\alpha_{12}y_2,
 \qquad
 z_2=\alpha_{21}y_1+\alpha_{22}y_2,
 \qquad
 A=|\alpha_{ik}|\ne0,
\]
et il est donc identique au plus petit commun multiple des deux modules premiers entre eux $\mN_1$ et $\mN_2$, qui possèdent respectivement les intégrales $z_1$ et $z_2$, et dont les polynômes générateurs sont donnés par
% unit: noether.p17.s12.eq.004
\[
 N_i=z_i\xi-z_i'=(\alpha_{i1}y_1+\alpha_{i2}y_2)\xi-(\alpha_{i1}y_1'+\alpha_{i2}y_2').
\]
Ainsi $\mM$ admet une infinité de représentations distinctes comme plus petit commun multiple de modules premiers entre eux, et, d'après le théorème VII, les deux modules $\mN_1$ et $\mN_2$ sont donc de même espèce entre eux et avec tous les $\mM_1$ et $\mM_2$. En fait, les groupes quotients de tous ces modules ont l'ordre 1, donc ils sont isomorphes.

% unit: noether.p17.s12.par.003
Nous voulons maintenant établir la décomposition additive du groupe quotient $\mG=\mA_1+\mA_2$ correspondant à la représentation (41). À cette fin, nous partons, comme au \S~4, de la relation (42), qui exprime que $\mM_1$ et $\mM_2$ sont premiers entre eux. Pour les unités $a_1$ et $a_2$, nous obtenons les classes résiduelles engendrées par les polynômes
% unit: noether.p17.s12.eq.005
\[
 -\frac{y_1}{D}(y_2\xi-y_2')
 \quad\text{et}\quad
 \frac{y_2}{D}(y_1\xi-y_1').
\]
Ainsi
% unit: noether.p17.s12.eq.006
\[
 \mM_1=\left(\mM,\frac{y_2}{D}(y_1\xi-y_1')\right),
 \qquad
 \mM_2=\left(\mM,-\frac{y_1}{D}(y_2\xi-y_2')\right).
\]
Si nous formons maintenant de même les unités $b_1$ et $b_2$ qui correspondent à la décomposition $\mM=[\mN_1,\mN_2]$, en remplaçant $y_i$ par $z_i$, il vient
% unit: noether.p17.s12.eq.007
\[
 -\frac{z_1}{D(z)}(z_2\xi-z_2')
 =-\frac{\alpha_{11}y_1+\alpha_{12}y_2}{AD(y)}
 ((\alpha_{21}y_1+\alpha_{22}y_2)\xi-(\alpha_{21}y_1'+\alpha_{22}y_2')),
\]
% unit: noether.p17.s12.eq.008
\[
 \frac{z_2}{D(z)}(z_1\xi-z_1')
 =\frac{\alpha_{21}y_1+\alpha_{22}y_2}{AD(y)}
 ((\alpha_{11}y_1+\alpha_{12}y_2)\xi-(\alpha_{11}y_1'+\alpha_{12}y_2')),
\]
donc
% unit: noether.p17.s12.eq.043
\begin{equation}
\left\{
\begin{aligned}
 b_1&=\frac{\alpha_{11}y_1+\alpha_{12}y_2}{y_1}
 \left(\frac{\alpha_{22}}{A}\right)a_1
 +\frac{\alpha_{11}y_1+\alpha_{12}y_2}{y_2}
 \left(-\frac{\alpha_{21}}{A}\right)a_2
 =\frac{z_1}{y_1}\alpha^{11}a_1+\frac{z_1}{y_2}\alpha^{12}a_2,\\
 b_2&=\frac{\alpha_{21}y_1+\alpha_{22}y_2}{y_1}
 \left(-\frac{\alpha_{12}}{A}\right)a_1
 +\frac{\alpha_{21}y_1+\alpha_{22}y_2}{y_2}
 \left(\frac{\alpha_{11}}{A}\right)a_2
 =\frac{z_2}{y_1}\alpha^{21}a_1+\frac{z_2}{y_2}\alpha^{22}a_2.
\end{aligned}\right.
\tag{43}
\end{equation}
Ici $(\alpha^{ik})$ est l'inverse de la matrice $(\alpha_{ik})$.

% unit: noether.p17.s12.eq.044
Inversement, en échangeant $z_i$ avec $y_i$, on obtient
\begin{equation}
\left\{
\begin{aligned}
 a_1&=\frac{y_1}{z_1}\alpha_{11}b_1+\frac{y_1}{z_2}\alpha_{12}b_2,\\
 a_2&=\frac{y_2}{z_1}\alpha_{21}b_1+\frac{y_2}{z_2}\alpha_{22}b_2.
\end{aligned}\right.
\tag{44}
\end{equation}
Ces formules font voir l'isomorphie des groupes $\mA$ et $\mB$. D'après nos considérations générales, en effet, si $a_1b_\sigma\ne0$, la correspondance en question est $a_1\sim a_1b_\sigma$; donc, d'après (44),
% unit: noether.p17.s12.eq.009
\[
 a_1\sim\frac{y_1}{z_\sigma}\alpha_{1\sigma}b_\sigma,
 \qquad\text{si }\alpha_{1\sigma}\ne0,
\]
et, d'après (43), de même
% unit: noether.p17.s12.eq.010
\[
 b_\sigma\sim\frac{z_\sigma}{y_2}\alpha^{\sigma2}a_2,
 \qquad\text{si }\alpha^{\sigma2}\ne0,
\]
d'où
% unit: noether.p17.s12.eq.011
\[
 \frac{y_1}{z_\sigma}\alpha_{1\sigma}b_\sigma
 \sim
 \frac{y_1}{z_\sigma}\alpha_{1\sigma}\frac{z_\sigma}{y_2}\alpha^{\sigma2}a_2
 =\frac{y_1}{y_2}\alpha_{1\sigma}\alpha^{\sigma2}a_2.
\]
Ainsi
% unit: noether.p17.s12.eq.012
\[
 a_1\sim\alpha_{1\sigma}\alpha^{\sigma2}\frac{y_1}{y_2}a_2
 \quad\text{pour }\alpha_{1\sigma}\alpha^{\sigma2}\ne0,
 \qquad
 \frac{1}{\alpha_{1\sigma}\alpha^{\sigma2}}\frac{y_2}{y_1}a_1\sim a_2,
\]
ce qui représente l'inverse. Si maintenant $(\alpha_{ik})$ n'est pas une matrice diagonale, cette condition $\alpha_{1\sigma}\alpha^{\sigma2}\ne0$ est toujours satisfaite pour un $\sigma$. Le cas de la matrice diagonale doit être exclu, car alors les modules sont identiques. Si, par exemple, seul $\alpha_{12}=0$, alors $\mA_1=\mB_1$, mais non $\mA_2=\mB_2$. De $\mA_1+\mA_2=\mB_1+\mB_2$ et $\mA_1=\mB_1$ on ne peut donc pas conclure que $\mA_2=\mB_2$ (cf. la remarque 12). Les trois modules $\mM_1$, $\mM_2$ et $\mN_2$ sont alors deux à deux premiers entre eux, mais ne forment pas un système totalement premier entre eux, puisque $[\mM_1,\mM_2]$ est divisible par $\mN_2$ (cf. la remarque 15).

% unit: noether.p17.s12.par.004
En posant encore
% unit: noether.p17.s12.eq.013
\[
 t_{12}=\alpha_{1\sigma}\alpha^{\sigma2}\frac{y_1}{y_2}a_2,
 \qquad
 t_{21}=\frac{1}{\alpha_{1\sigma}\alpha^{\sigma2}}\frac{y_2}{y_1}a_1,
\]
les relations (35), comme on le confirme par le calcul, sont satisfaites.

% unit: noether.p17.s12.par.005
2. Un exemple de groupe infini est fourni par tout module à plus d'une variable dont la base ne se compose que d'un seul polynôme. L'exemple suivant, un peu plus compliqué, doit montrer que des groupes infinis peuvent aussi apparaître pour des modules à plusieurs éléments, c'est-à-dire pour ceux dont la base du module comprend plus d'un polynôme.\footnote{Comparer l'exemple de Landau communiqué par Blumberg (loc. cit., p. 51--52). Ici et dans la suite nous désignons les variables indépendantes par $x$ et $y$.}
% unit: noether.p17.s12.eq.014
\[
 \mM=[(\xi+x\eta),(\xi+1)]=[(P),(Q)]
 \qquad\left(\xi=\frac{\partial}{\partial x},\ \eta=\frac{\partial}{\partial y}\right).
\]
La règle de multiplication est celle des expressions différentielles. Le groupe quotient est ici infini parce que les groupes quotients de $(\xi+x\eta)$ et de $(\xi+1)$ sont tous deux infinis, tandis que les modules $(\xi+x\eta)$ et $(\xi+1)$ sont premiers entre eux, car
% unit: noether.p17.s12.eq.015
\[
 1=(-x\xi-x^2\eta+1)Q+(x\xi+x-1)P.
\]
Il est maintenant très facile de montrer que $\mM$ ne possède aucun polynôme du second degré, en posant
% unit: noether.p17.s12.eq.016
\[
 (u_1\xi+u_2\eta+u_3)(\xi+x\eta)=(v_1\xi+v_2\eta+v_3)(\xi+1).
\]
On en tire $u_1=u_2=u_3=v_1=v_2=v_3=0$. Si l'on cherche ensuite des polynômes du troisième degré dans $\mM$, on en trouve deux linéairement indépendants, par exemple
% unit: noether.p17.s12.eq.017
\[
 (\xi^2+x\xi\eta+\xi+(2+x)\eta)Q=(\xi^2+2\xi+1)P
\]
et
% unit: noether.p17.s12.eq.018
\[
 (\xi^2+2x\xi\eta+x^2\eta^2+\eta)Q
 =(\xi^2+x\xi\eta+\xi+(x-1)\eta)P.
\]
Si $\mM$ était un module à un seul élément, les deux polynômes devraient être divisibles par le polynôme générateur et, puisque $\mM$ ne contient aucun polynôme du second degré, par un polynôme du troisième degré; ils devraient donc être proportionnels, ce qui n'est pas le cas. Le fait que $\mM$ ait plusieurs éléments est la raison interne de l'échec des théorèmes de produit de Landau, valables dans le cas d'une seule variable.

% unit: noether.p17.s12.par.006
3. Un autre exemple montre que des groupes infinis peuvent aussi apparaître pour des modules premiers.\footnote{Notre définition de \og{} module premier \fg{} est différente de ce que l'on appelle module premier dans le cas commutatif, car il existe alors toujours des diviseurs de multiplicité inférieure, sauf lorsque la multiplicité est $0$, c'est-à-dire lorsque le groupe quotient est fini.} Soit $P$ le domaine de toutes les fonctions rationnelles de deux variables $x,y$. Le module $\mM$ de base
% unit: noether.p17.s12.eq.019
\[
 \xi-\frac{y}{x}\eta=\xi-a
\]
possède un groupe infini, puisque celui-ci contient la classe résiduelle de chaque puissance de $y$; d'autre part, $\mM$ est un module premier. En effet, nous montrons que tout diviseur
% unit: noether.p17.s12.eq.020
\[
 \mM_1=(\xi-a,F(\xi,\eta))
\]
est égal au module unité. Tout polynôme $F(\xi,\eta)$ peut en effet être représenté modulo $(\xi-a)$ par un polynôme $G(\eta)$:
% unit: noether.p17.s12.eq.021
\[
 F(\xi,\eta)\equiv c_0\eta^\nu+c_1\eta^{\nu-1}+\cdots+c_\nu\quad(\mM).
\]
Sans restreindre la généralité, soit $c_0=1$. Or, d'après la loi de multiplication, $\xi\eta^\nu-\eta^\nu\xi=0$; donc, puisque
% unit: noether.p17.s12.eq.022
\[
 \xi\equiv a(\mM_1),
 \qquad
 \eta^\nu\equiv F_1(\mM_1),
\]
où $F_1=-(c_1\eta^{\nu-1}+\cdots+c_\nu)$, on a
% unit: noether.p17.s12.eq.023
\[
 \xi F_1-\eta^\nu a\equiv0(\mM_1).
\]
En développant
% unit: noether.p17.s12.eq.024
\[
 \eta^\nu a=a\eta^\nu+\nu\frac{\partial a}{\partial y}\eta^{\nu-1}+\cdots,
\]
et en remplaçant de nouveau ici $\eta^\nu$ par $F_1$, on obtient
% unit: noether.p17.s12.eq.025
\[
 \xi(c_1\eta^{\nu-1}+\cdots+c_\nu)-a(c_1\eta^{\nu-1}+\cdots+c_\nu)
 +\nu\frac{\partial a}{\partial y}\eta^{\nu-1}+\cdots\equiv0(\mM_1),
\]
ou
% unit: noether.p17.s12.eq.026
\[
 c_1\eta^{\nu-1}a+\frac{\partial c_1}{\partial x}\eta^{\nu-1}+\cdots
 -ac_1\eta^{\nu-1}+\cdots+
 \nu\frac{\partial a}{\partial y}\eta^{\nu-1}+\cdots
 \equiv0(\mM_1),
\]
ou enfin
% unit: noether.p17.s12.eq.027
\[
 \left(\frac{\partial c_1}{\partial x}+\nu\frac{\partial a}{\partial y}\right)\eta^{\nu-1}+\cdots
 \equiv0(\mM_1).
\]
À gauche se trouve maintenant un polynôme en $\eta$ de degré exact $\nu-1$, qui ne s'annule pas identiquement; car son coefficient dominant
% unit: noether.p17.s12.eq.028
\[
 \frac{\partial c_1}{\partial x}+\nu\frac{\partial a}{\partial y}
 =\frac{\partial c_1}{\partial x}+\nu\frac{1}{x}
 \quad\text{est }\ne0,
\]
puisque sinon $c_1=-\nu\log x+\varphi(y)$ ne serait pas rationnel. Le procédé peut donc être poursuivi et conduit à un polynôme non identiquement nul de degré zéro, ce qui montre que l'unité est aussi contenue dans $\mM_1$.\footnote{Ce procédé revient à démontrer que les conditions d'intégrabilité nécessaires pour un système d'équations différentielles ne sont pas satisfaites.}

% unit: noether.p17.s12.par.007
4. Si, en revanche, nous adjoignons encore au domaine de rationalité la fonction $\log x$, alors $\mM$ possède les diviseurs
% unit: noether.p17.s12.eq.029
\[
 \left(\xi-\frac{y}{x},\ \eta-\log x+\varphi(y)\right),
\]
où $\varphi(y)$ parcourt toutes les fonctions rationnelles de $y$. Ces diviseurs sont tous distincts les uns des autres; pour chacun la condition d'intégrabilité est satisfaite, et l'intégrale correspondante est
% unit: noether.p17.s12.eq.030
\[
 y\log x-\int\varphi(y)\,dy+c.
\]
Par conséquent, l'unité ne peut pas être contenue dans le module, car sinon la fonction nulle serait la seule intégrale possible. D'autre part, $\xi$ et $\eta$ sont, relativement au module, congrues à une grandeur du domaine; le groupe quotient est donc fini et d'ordre 1. De plus, deux quelconques de ces diviseurs en nombre infini sont premiers entre eux, et le module donné $\mM$ est égal au plus petit commun multiple $\mN$ de tous ces diviseurs. En effet, il est divisible par tous les diviseurs, donc aussi par $\mN$. Si maintenant $\mN$ était un diviseur propre de $\mM$, alors $\mN$ devrait encore contenir au moins un polynôme $F$, disons de degré $\rho$, qui ne dépende que de $\eta$. Faisons parcourir à $\varphi(y)$ les fonctions $y,\ldots,y^{\rho+1}$; il n'existe alors aucune relation linéaire, à coefficients indépendants de $y$, entre les intégrales correspondantes
% unit: noether.p17.s12.eq.031
\[
 y\log x-\frac{y^{i+1}}{i+1}+c_i\qquad(i=1,\ldots,\rho+1).
\]
Mais l'équation différentielle d'ordre $\rho$, $F=0$, ne peut pas posséder $\rho+1$ intégrales linéairement indépendantes.

% unit: noether.p17.s12.par.008
Le module $\mM=\mN$ n'est cependant pas complètement réductible.\footnote{Nous avons ainsi un exemple du premier cas du théorème X.} Sinon, il serait le plus petit commun multiple d'un nombre fini de ces modules premiers, et son groupe quotient serait donc d'ordre fini, tandis que le groupe quotient de $\xi-y/x$ contient bien les classes résiduelles de toutes les puissances de $\eta$, et est donc infini.

% unit: noether.p17.close.001
Göttingen, le 1er août 1919.

\begin{center}
(Reçu le 4 août 1919.)
\end{center}
